\markdownRendererDocumentBegin
\markdownRendererSectionBegin
\markdownRendererSectionBegin
\markdownRendererHeadingTwo{Patrones y Diseño de Algoritmos}\markdownRendererInterblockSeparator
{}\markdownRendererSectionBegin
\markdownRendererHeadingThree{Hamming Distance}\markdownRendererParagraphSeparator
{}La distancia de Hamming entre dos cadenas \markdownRendererDollarSign{}a\markdownRendererDollarSign{} y \markdownRendererDollarSign{}b\markdownRendererDollarSign{} de igual longitud es el número de posiciones donde difieren. Por ejemplo, \markdownRendererDollarSign{}hamming(01101, 11001) = 2\markdownRendererDollarSign{}.\markdownRendererParagraphSeparator
{}Si representamos las cadenas como enteros (cuando \markdownRendererDollarSign{}k \markdownRendererBackslash{}le 32\markdownRendererDollarSign{}), podemos calcular la distancia usando operaciones a nivel de bit. El XOR (\markdownRendererDollarSign{}\markdownRendererBackslash{}oplus\markdownRendererDollarSign{}) identifica los bits diferentes y contar los 1s.\markdownRendererParagraphSeparator
{}Esto reduce la complejidad de O(\markdownRendererDollarSign{}n\markdownRendererCircumflex{}2 k\markdownRendererDollarSign{}) a O(\markdownRendererDollarSign{}n\markdownRendererCircumflex{}2\markdownRendererDollarSign{}) al comparar \markdownRendererDollarSign{}n\markdownRendererDollarSign{} cadenas de longitud \markdownRendererDollarSign{}k\markdownRendererDollarSign{}.\markdownRendererInterblockSeparator
{}\markdownRendererInputFencedCode{./_markdown_main/8582959398df3da9526e553164206b82.verbatim}{cpp}{cpp}\markdownRendererInterblockSeparator
{}\markdownRendererDollarSign{}\markdownRendererBackslash{}clearpage\markdownRendererDollarSign{}\markdownRendererInterblockSeparator
{}
\markdownRendererSectionEnd \markdownRendererSectionBegin
\markdownRendererHeadingThree{Reachability in Graphs}\markdownRendererInterblockSeparator
{}Dado un DAG con \markdownRendererDollarSign{}n\markdownRendererDollarSign{} nodos, queremos calcular \markdownRendererDollarSign{}reach(x)\markdownRendererDollarSign{}: número de nodos alcanzables desde \markdownRendererDollarSign{}x\markdownRendererDollarSign{}.\markdownRendererParagraphSeparator
{}Se puede usar programación dinámica clásica O(\markdownRendererDollarSign{}n\markdownRendererCircumflex{}2\markdownRendererDollarSign{}) o bit-paralela usando \markdownRendererCodeSpan{bitset} para cada nodo, donde \markdownRendererCodeSpan{reach[x] \markdownRendererPipe{}= reach[u]} para cada \markdownRendererDollarSign{}u\markdownRendererDollarSign{} adyacente a \markdownRendererDollarSign{}x\markdownRendererDollarSign{}. Esto reduce operaciones por nodo usando OR de bits, pero puede usar mucha memoria (\markdownRendererDollarSign{}O(n\markdownRendererCircumflex{}2)\markdownRendererDollarSign{} bits).\markdownRendererInterblockSeparator
{}\markdownRendererInputFencedCode{./_markdown_main/101b1ad853a3fc625f34f6f262366a71.verbatim}{cpp}{cpp}\markdownRendererInterblockSeparator
{}
\markdownRendererSectionEnd \markdownRendererSectionBegin
\markdownRendererHeadingThree{Counting Subgrids}\markdownRendererInterblockSeparator
{}Dado un grid \markdownRendererDollarSign{}n \markdownRendererBackslash{}times n\markdownRendererDollarSign{} con casillas negras (1) o blancas (0), queremos contar subgrids cuyos cuatro vértices son negros.\markdownRendererParagraphSeparator
{}En la implementación bit-paralela, cada fila se representa como un \markdownRendererCodeSpan{bitset} y usamos operaciones AND para contar columnas donde ambas filas tienen negro. La complejidad baja de O(\markdownRendererDollarSign{}n\markdownRendererCircumflex{}3\markdownRendererDollarSign{}) a O(\markdownRendererDollarSign{}n\markdownRendererCircumflex{}2\markdownRendererDollarSign{}) y es mucho más rápida, pero usa más memoria (\markdownRendererDollarSign{}O(n\markdownRendererCircumflex{}2)\markdownRendererDollarSign{} bits para los bitsets).\markdownRendererInterblockSeparator
{}\markdownRendererInputFencedCode{./_markdown_main/efd015706bd2c9c5f243ce93c03cf8fb.verbatim}{cpp}{cpp}\markdownRendererInterblockSeparator
{}\markdownRendererDollarSign{}\markdownRendererBackslash{}clearpage\markdownRendererDollarSign{}\markdownRendererInterblockSeparator
{}
\markdownRendererSectionEnd \markdownRendererSectionBegin
\markdownRendererHeadingThree{Two Pointers Method}\markdownRendererInterblockSeparator
{}Se usan dos punteros \markdownRendererDollarSign{}l\markdownRendererDollarSign{} y \markdownRendererDollarSign{}r\markdownRendererDollarSign{} que solo avanzan hacia la derecha para encontrar subarrays con ciertas propiedades (ej: suma = x).\markdownRendererParagraphSeparator
{}Ejemplo clásico: encontrar subarray con suma exacta x. Ambos punteros se mueven O(\markdownRendererDollarSign{}n\markdownRendererDollarSign{}) pasos en total, dando complejidad O(\markdownRendererDollarSign{}n\markdownRendererDollarSign{}).\markdownRendererParagraphSeparator
{}También se aplica al problema 2SUM: encontrar dos valores cuya suma sea x. Primero se ordena el array O(\markdownRendererDollarSign{}n\markdownRendererBackslash{}log n\markdownRendererDollarSign{}) y luego se usan dos punteros desde extremos opuestos.\markdownRendererInterblockSeparator
{}\markdownRendererInputFencedCode{./_markdown_main/509de053210700908a0ccb310219c5aa.verbatim}{cpp}{cpp}\markdownRendererInterblockSeparator
{}
\markdownRendererSectionEnd \markdownRendererSectionBegin
\markdownRendererHeadingThree{Nearest Smaller Element (Amortized Stack)}\markdownRendererInterblockSeparator
{}Para cada elemento de un array, encontramos el primer elemento más pequeño que lo precede usando un stack.\markdownRendererParagraphSeparator
{}Cada elemento se agrega exactamente una vez y se elimina a lo sumo una vez → O(\markdownRendererDollarSign{}n\markdownRendererDollarSign{}) tiempo. Es un ejemplo de análisis amortizado: operaciones individuales varían, pero el total es lineal.\markdownRendererInterblockSeparator
{}\markdownRendererInputFencedCode{./_markdown_main/1a86719234dc2cbc9f1ec51d05d60072.verbatim}{cpp}{cpp}\markdownRendererInterblockSeparator
{}
\markdownRendererSectionEnd \markdownRendererSectionBegin
\markdownRendererHeadingThree{Sliding Window Minimum}\markdownRendererInterblockSeparator
{}Se mantiene una ventana de tamaño fijo moviéndose sobre un array, y queremos el mínimo de cada ventana.\markdownRendererParagraphSeparator
{}Se usa un deque: los elementos se agregan al final y se eliminan del final si son mayores al nuevo elemento. El primero siempre es el mínimo. Cada elemento se agrega y elimina a lo sumo una vez → O(\markdownRendererDollarSign{}n\markdownRendererDollarSign{}).\markdownRendererParagraphSeparator
{}\markdownRendererInputFencedCode{./_markdown_main/4edc1587ad92a9bbbf6d818bef3b7e4e.verbatim}{cpp}{cpp}\markdownRendererInterblockSeparator
{}
\markdownRendererSectionEnd \markdownRendererSectionBegin
\markdownRendererHeadingThree{Minimizing Sums}\markdownRendererInterblockSeparator
{}Dado un conjunto de \markdownRendererDollarSign{}n\markdownRendererDollarSign{} números \markdownRendererDollarSign{}a\markdownRendererUnderscore{}1, a\markdownRendererUnderscore{}2, \markdownRendererBackslash{}dots, a\markdownRendererUnderscore{}n\markdownRendererDollarSign{}, consideremos el problema de encontrar un valor \markdownRendererDollarSign{}x\markdownRendererDollarSign{} que minimice la suma\markdownRendererParagraphSeparator
{}\markdownRendererDollarSign{}\markdownRendererDollarSign{}\markdownRendererPipe{}a\markdownRendererUnderscore{}1 - x\markdownRendererPipe{} + \markdownRendererPipe{}a\markdownRendererUnderscore{}2 - x\markdownRendererPipe{} + \markdownRendererBackslash{}dots + \markdownRendererPipe{}a\markdownRendererUnderscore{}n - x\markdownRendererPipe{}\markdownRendererDollarSign{}\markdownRendererDollarSign{}\markdownRendererParagraphSeparator
{}Por ejemplo, si los números son [1, 2, 9, 2, 6], la solución óptima es \markdownRendererDollarSign{}x = 2\markdownRendererDollarSign{}, lo que produce la suma\markdownRendererParagraphSeparator
{}\markdownRendererDollarSign{}\markdownRendererDollarSign{}\markdownRendererPipe{}1-2\markdownRendererPipe{} + \markdownRendererPipe{}2-2\markdownRendererPipe{} + \markdownRendererPipe{}9-2\markdownRendererPipe{} + \markdownRendererPipe{}2-2\markdownRendererPipe{} + \markdownRendererPipe{}6-2\markdownRendererPipe{} = 12\markdownRendererDollarSign{}\markdownRendererDollarSign{}\markdownRendererParagraphSeparator
{}Cada función \markdownRendererDollarSign{}\markdownRendererPipe{}a\markdownRendererUnderscore{}k - x\markdownRendererPipe{}\markdownRendererDollarSign{} es convexa, así que la suma también lo es. Podríamos usar búsqueda ternaria para encontrar el \markdownRendererDollarSign{}x\markdownRendererDollarSign{} óptimo, pero hay una solución más simple: el valor óptimo de \markdownRendererDollarSign{}x\markdownRendererDollarSign{} siempre es la mediana de los números, es decir, el elemento central después de ordenar. Por ejemplo, [1, 2, 9, 2, 6] se ordena como [1, 2, 2, 6, 9], y la mediana es 2.\markdownRendererParagraphSeparator
{}La mediana siempre es óptima, porque si \markdownRendererDollarSign{}x\markdownRendererDollarSign{} es menor que la mediana, la suma disminuye al aumentar \markdownRendererDollarSign{}x\markdownRendererDollarSign{}, y si \markdownRendererDollarSign{}x\markdownRendererDollarSign{} es mayor, la suma disminuye al disminuir \markdownRendererDollarSign{}x\markdownRendererDollarSign{}. Si \markdownRendererDollarSign{}n\markdownRendererDollarSign{} es par y hay dos medianas, ambas medianas y todos los valores entre ellas son soluciones óptimas.\markdownRendererParagraphSeparator
{}Ahora consideremos minimizar la función\markdownRendererParagraphSeparator
{}\markdownRendererDollarSign{}\markdownRendererDollarSign{}(a\markdownRendererUnderscore{}1 - x)\markdownRendererCircumflex{}2 + (a\markdownRendererUnderscore{}2 - x)\markdownRendererCircumflex{}2 + \markdownRendererBackslash{}dots + (a\markdownRendererUnderscore{}n - x)\markdownRendererCircumflex{}2\markdownRendererDollarSign{}\markdownRendererDollarSign{}\markdownRendererParagraphSeparator
{}Por ejemplo, con [1, 2, 9, 2, 6], la mejor elección es \markdownRendererDollarSign{}x = 4\markdownRendererDollarSign{}, produciendo la suma\markdownRendererParagraphSeparator
{}\markdownRendererDollarSign{}\markdownRendererDollarSign{}(1-4)\markdownRendererCircumflex{}2 + (2-4)\markdownRendererCircumflex{}2 + (9-4)\markdownRendererCircumflex{}2 + (2-4)\markdownRendererCircumflex{}2 + (6-4)\markdownRendererCircumflex{}2 = 46\markdownRendererDollarSign{}\markdownRendererDollarSign{}\markdownRendererParagraphSeparator
{}Esta función también es convexa. Aunque podríamos usar búsqueda ternaria, hay una solución directa: el valor óptimo de \markdownRendererDollarSign{}x\markdownRendererDollarSign{} es el promedio de los números. En este ejemplo, el promedio es\markdownRendererParagraphSeparator
{}\markdownRendererDollarSign{}\markdownRendererDollarSign{}(1 + 2 + 9 + 2 + 6)/5 = 4\markdownRendererDollarSign{}\markdownRendererDollarSign{}\markdownRendererParagraphSeparator
{}Esto se puede demostrar reescribiendo la suma como\markdownRendererParagraphSeparator
{}\markdownRendererDollarSign{}\markdownRendererDollarSign{}n x\markdownRendererCircumflex{}2 - 2x(a\markdownRendererUnderscore{}1 + a\markdownRendererUnderscore{}2 + \markdownRendererBackslash{}dots + a\markdownRendererUnderscore{}n) + (a\markdownRendererUnderscore{}1\markdownRendererCircumflex{}2 + a\markdownRendererUnderscore{}2\markdownRendererCircumflex{}2 + \markdownRendererBackslash{}dots + a\markdownRendererUnderscore{}n\markdownRendererCircumflex{}2)\markdownRendererDollarSign{}\markdownRendererDollarSign{}\markdownRendererParagraphSeparator
{}La última parte no depende de \markdownRendererDollarSign{}x\markdownRendererDollarSign{}, y la función restante forma una parábola con mínimo en \markdownRendererDollarSign{}x = s/n\markdownRendererDollarSign{}, donde \markdownRendererDollarSign{}s = a\markdownRendererUnderscore{}1 + a\markdownRendererUnderscore{}2 + \markdownRendererBackslash{}dots + a\markdownRendererUnderscore{}n\markdownRendererDollarSign{}. Por lo tanto, el mínimo se alcanza en el promedio de los números \markdownRendererDollarSign{}a\markdownRendererUnderscore{}1, a\markdownRendererUnderscore{}2, \markdownRendererBackslash{}dots, a\markdownRendererUnderscore{}n\markdownRendererDollarSign{}.\markdownRendererInterblockSeparator
{}\markdownRendererInputFencedCode{./_markdown_main/61b6e729390ee22d3ad135987fbffa2e.verbatim}{cpp}{cpp}
\markdownRendererSectionEnd 
\markdownRendererSectionEnd 
\markdownRendererSectionEnd \markdownRendererDocumentEnd